% User-facing changes in the v1.10 source tree. Keep before the inherited chapters.
\clearpage
\section{What changed in v1.10}\label{sec:v110-user}

The chapters following this section were copied from the v1.00 manual. They
describe the physical options and output schemas of that version. Their
planet-specific timings, residuals, and convergence examples are historical
v1.00 measurements; they are not renewed v1.10 results. This section gives the
current v1.10 behavior where the versions differ. In particular, the older
use of the word ``converged'' for a marching snapshot does not certify a
stationary solution.

\subsection{Two levels of convergence}

\textbf{Regional stabilization} is the finite-duration test implemented for
atomic global physical-time marching. It checks whether upper-atmosphere
profiles have small variations while lower-atmosphere variations remain within
the configured bounds. It ends a successful regional run without launching a
Newton--Krylov solve. \textbf{Stationary convergence} requires the enabled
fluid, chemical, radiative, and boundary equations to satisfy specified
stationary residual and consistency tolerances. Regional stabilization alone
does not imply a stationary root. The numerical method does not define either
physical criterion.

Enable the regional test with \texttt{Regional convergence: True}. The main
program accepts it only with \texttt{Run mode: phys} in an atomic calculation
using global physical time steps, with Newton solving disabled. A local-time
step, molecular chemistry, stationary restart, or post-processing-only run with
this option stops at startup. The absence of the key leaves the regional
monitor disabled. The Huang and Salz settings under
\texttt{examples/regional\_convergence/} are fragments; merge one into a
complete planetary \texttt{input.inp}, replacing duplicate keys.

\begin{center}
\begin{tabular}{lcc}
\toprule
Quantity checked & Upper region & Lower/transition regions\\
\midrule
Window-mean change and span & 0.01 & 0.10\\
Full value range & 0.05 & 0.20\\
Mass and internal-energy inventory & 0.01 & 0.05\\
Mean actual face mass flow & 0.01 & 0.05\\
\bottomrule
\end{tabular}
\end{center}

These are source defaults, not universal observational accuracy bounds. The
three configured regions use their own limits. The inventory limit applies to
both the change and the full range. Mass flow is checked through its window
means and their span; its instantaneous range is recorded but is not a stop
gate. All criteria must hold over four validation windows times a confirmation
factor of two, for eight consecutive windows, with at least 40 accepted
endpoints in each. Rejected time steps do not advance the confirmation clock.
When a candidate fails, the test starts a new candidate with a new fixed
reference; it does not combine favorable windows from different candidates.

The radial cuts, window duration, settling time, total time budget, thresholds,
and profile cadence are editable with \texttt{Regional ...} input lines. The
complete v1.10 set is:

\begin{lstlisting}[style=cfg]
Regional convergence: True
Regional lower radius: 1.3
Regional upper radius: 1.85
Regional window seconds: 87965.04804471003
Regional maximum seconds: 0
Regional settling seconds: 0
Regional validation windows: 4
Regional confirmation factor: 2
Regional minimum samples: 40
Regional lower mean limit: 0.10
Regional transition mean limit: 0.10
Regional upper mean limit: 0.01
Regional lower range limit: 0.20
Regional transition range limit: 0.20
Regional upper range limit: 0.05
Regional lower inventory limit: 0.05
Regional transition inventory limit: 0.05
Regional upper inventory limit: 0.01
Regional lower mass flow limit: 0.05
Regional transition mass flow limit: 0.05
Regional upper mass flow limit: 0.01
Regional profile interval: 40
\end{lstlisting}

This fragment shows the Huang radial cuts and window; it is not a complete
input file. A common \texttt{Regional inventory limit} or
\texttt{Regional mass flow limit} may replace the three corresponding regional
limits. An explicit regional value takes precedence over the common value
regardless of line order. Unknown \texttt{Regional ...} labels stop the run.
If \texttt{Regional window seconds: 0}, the duration is twice the initial
outward acoustic travel time. If \texttt{Regional maximum seconds: 0}, the
budget is the settling time plus two complete confirmation intervals; zero
does not mean unlimited time. The initial lower and upper cuts default to
$1.2\,R_{\rm p}$ when no cuts are supplied, so an empty transition region is
allowed. The cuts and duration should be stated when comparing calculations.

\subsection{Progress records and plotting}

The regional monitor writes three files:
\begin{itemize}
\item \path{output/regional_convergence.csv}: criterion history;
\item \path{output/regional_convergence_status.txt}: effective settings and
current outcome;
\item \path{output/regional_profile_history.txt}: profiles with region labels.
\end{itemize}
The monitor evaluates every accepted step. Profiles and the collecting status
are updated every 40 accepted steps by default, with initial and terminal
profile records. The status file reports
\path{convergence_level=regional_stabilization} only after a complete
successful candidate and \path{stationary_root_claim=False} in that case.
An incomplete run reports \path{convergence_level=none}. When regional
monitoring is enabled but the specified criteria are not met by normal
finalization, the program exits with status 4; a separate startup or runtime
failure can end earlier.

Run the plotter from the calculation directory after output appears:

\begin{lstlisting}[style=sh]
python3 /path/to/EXHALE_v1.10/EXHALE_plots.py --regional --live 60
python3 /path/to/EXHALE_v1.10/EXHALE_plots.py --regional --save output/regional
\end{lstlisting}

The first invocation refreshes the display every 60 seconds. The second
writes PDF figures. The regional view shows physical profiles and progress
separately for the lower, transition, and upper regions. The plotter reads
existing output; it does not advance a calculation. The old profile view is
available with \texttt{--profiles}.

\subsection{Stationary lower-boundary closure and current scope}

The v1.10 stationary residual closes the lower ghost composition and
thermodynamics at a fixed column. Certification of the stored state also
checks the boundary result; an unclosed boundary cannot establish a certified
stationary state. The relevant modules are
\texttt{src/modules/states/stationary\_boundary\_chemistry.f90} and
\texttt{reservoir\_composition.f90}. This implementation does not by itself
establish stationary convergence for the historical Yan, Huang, or Salz
states. No such claim follows from the new regional test.

The older \texttt{Max steps} input and \path{EXHALE_MAXSTEPS}
environment control set separate caps. Similarly,
\path{EXHALE_MARCH_FLUX_TH} changes the marching stop, while
\texttt{Flux spread tol} belongs to stationary acceptance. These values must
not be substituted for one another. \path{EXHALE_MONITOR_STEPS} must be a
positive integer in the current source. The copied v1.00 descriptions of
experimental environment settings do not establish that they are documented
user inputs in v1.10.

The v1.10 tree lacks the v1.00 regression harness. Its \texttt{make check}
target reports the missing harness as a failure. The README in this tree
describes source-build requirements and the status of the paused calculation
sequence. No model run was made to produce this documentation update.
